\section{Stage 2: From the PDE to CN/PSOR}

\subsection{Purpose}
This stage explains the finite-difference matrix construction in enough detail
to connect every mathematical coefficient with the arrays used by the solver.
The target is a reader who understands derivatives and matrix multiplication
but has not previously implemented an American-option LCP.

\subsection{Method}
\subsubsection{Spot and time grids}
Truncate the unbounded spot domain to \([0,S_{\max}]\). Choose \(M\) equal spot
intervals and \(N\) equal time intervals:
\[
  S_i=i\Delta S,
  \qquad
  i=0,\ldots,M,
  \qquad
  \Delta S=\frac{S_{\max}}{M},
\]
\[
  \tau_n=n\Delta\tau,
  \qquad
  n=0,\ldots,N,
  \qquad
  \Delta\tau=\frac{T}{N}.
\]
Write \(U_i^n\) for the approximation to \(U(S_i,\tau_n)\). At a fixed time,
the interior unknown vector is
\[
  U^n=(U_1^n,U_2^n,\ldots,U_{M-1}^n)^\top.
\]

\begin{commonconfusion}
The index \(i\) moves across spot and the index \(n\) moves across time to
maturity. \(U_i^n\) is one scalar. \(U^n\) is the column vector containing all
interior spot values at time level \(n\).
\end{commonconfusion}

\begin{commonconfusion}
\(U_0^n\) and \(U_M^n\) are not missing unknowns. They are values at the two
domain boundaries \(S=0\) and \(S=S_{\max}\), where an option-specific boundary
condition supplies the value.
\end{commonconfusion}

\subsection{Mathematical Formulation}
\stepheading{central differences}
At an interior node, central differences give
\[
  U_S(S_i,\tau_n)
  \approx
  \frac{U_{i+1}^n-U_{i-1}^n}{2\Delta S},
\]
\[
  U_{SS}(S_i,\tau_n)
  \approx
  \frac{U_{i+1}^n-2U_i^n+U_{i-1}^n}{(\Delta S)^2}.
\]
Substitute them into
\[
  \mathcal{L}U
  =
  \frac12\sigma^2S^2U_{SS}
  +(r-q)SU_S-rU.
\]
At \(S_i\),
\begin{align*}
  (L_hU)_i
  &=
  \frac12\sigma^2S_i^2
  \frac{U_{i+1}-2U_i+U_{i-1}}{(\Delta S)^2}\\
  &\quad
  +(r-q)S_i
  \frac{U_{i+1}-U_{i-1}}{2\Delta S}
  -rU_i.
\end{align*}
Collect the coefficients of the three neighboring values:
\[
  (L_hU)_i=a_iU_{i-1}+b_iU_i+c_iU_{i+1},
\]
where
\[
  a_i=
  \frac{\tfrac12\sigma^2S_i^2}{(\Delta S)^2}
  -\frac{(r-q)S_i}{2\Delta S},
\]
\[
  b_i=
  -\frac{\sigma^2S_i^2}{(\Delta S)^2}-r,
\]
\[
  c_i=
  \frac{\tfrac12\sigma^2S_i^2}{(\Delta S)^2}
  +\frac{(r-q)S_i}{2\Delta S}.
\]

\subsubsection{Where the coefficients come from}
\begin{center}
\begin{tabular}{>{\raggedright\arraybackslash}p{0.18\textwidth}
                >{\raggedright\arraybackslash}p{0.23\textwidth}
                >{\raggedright\arraybackslash}p{0.23\textwidth}
                >{\raggedright\arraybackslash}p{0.23\textwidth}}
\toprule
Source & Contribution to \(a_i\) & Contribution to \(b_i\) & Contribution to \(c_i\)\\
\midrule
Diffusion
& \(+\tfrac12\sigma^2S_i^2/(\Delta S)^2\)
& \(-\sigma^2S_i^2/(\Delta S)^2\)
& \(+\tfrac12\sigma^2S_i^2/(\Delta S)^2\)\\
Drift
& \(-(r-q)S_i/(2\Delta S)\)
& \(0\)
& \(+(r-q)S_i/(2\Delta S)\)\\
Discount
& \(0\)
& \(-r\)
& \(0\)\\
\bottomrule
\end{tabular}
\end{center}

\begin{readerintuition}
Diffusion spreads weight symmetrically to the left and right. Drift shifts the
two off-diagonal weights in opposite directions. Discounting affects only the
current node and therefore appears only in the diagonal coefficient.
\end{readerintuition}

\stepheading{a small matrix with \(M=4\)}
With nodes \(S_0,S_1,S_2,S_3,S_4\), there are three interior unknowns:
\[
  U^n=(U_1^n,U_2^n,U_3^n)^\top.
\]
The interior matrix is
\[
  L_h=
  \begin{bmatrix}
    b_1 & c_1 & 0\\
    a_2 & b_2 & c_2\\
    0 & a_3 & b_3
  \end{bmatrix}.
\]
The full stencil equations are
\begin{align*}
  (L_hU)_1&=a_1U_0+b_1U_1+c_1U_2,\\
  (L_hU)_2&=a_2U_1+b_2U_2+c_2U_3,\\
  (L_hU)_3&=a_3U_2+b_3U_3+c_3U_4.
\end{align*}
Matrix multiplication by the displayed \(3\times3\) matrix contains only the
interior terms. The omitted \(a_1U_0\) and \(c_3U_4\) are known boundary
contributions.

\begin{commonconfusion}
The first matrix row has no column for \(U_0\), and the last row has no column
for \(U_4\), because the matrix maps the three-component interior vector to
another three-component vector. Only the first stencil touches the lower
boundary; only the last stencil touches the upper boundary.
\end{commonconfusion}

\stepheading{Crank--Nicolson time stepping}
Crank--Nicolson averages the spatial operator at the old and new time levels:
\[
  \frac{U^{n+1}-U^n}{\Delta\tau}
  =
  \frac12L_hU^{n+1}
  +\frac12L_hU^n.
\]
Multiply by \(\Delta\tau\):
\[
  U^{n+1}-U^n
  =
  \frac{\Delta\tau}{2}L_hU^{n+1}
  +\frac{\Delta\tau}{2}L_hU^n.
\]
Move the new-time operator to the left:
\[
  U^{n+1}
  -\frac{\Delta\tau}{2}L_hU^{n+1}
  =
  U^n
  +\frac{\Delta\tau}{2}L_hU^n.
\]
Define
\[
  A=I-\frac{\Delta\tau}{2}L_h,
  \qquad
  B=I+\frac{\Delta\tau}{2}L_h.
\]
Then
\[
  AU^{n+1}=BU^n+g^{n,n+1},
\]
where \(g^{n,n+1}\) collects known boundary values.

\begin{commonconfusion}
\(A\) multiplies \(U^{n+1}\) because its terms were gathered on the left.
\(B\) multiplies \(U^n\) because its terms remain on the right. The signs are a
consequence of rearrangement, not an arbitrary naming convention.
\end{commonconfusion}

\subsubsection{What one row means}
Ignoring the boundary adjustment momentarily, an interior row \(i\) reads
\begin{align*}
 &-\frac{\Delta\tau}{2}a_iU_{i-1}^{n+1}
 +\left(1-\frac{\Delta\tau}{2}b_i\right)U_i^{n+1}
 -\frac{\Delta\tau}{2}c_iU_{i+1}^{n+1}\\
 &\qquad=
 \frac{\Delta\tau}{2}a_iU_{i-1}^{n}
 +\left(1+\frac{\Delta\tau}{2}b_i\right)U_i^{n}
 +\frac{\Delta\tau}{2}c_iU_{i+1}^{n}.
\end{align*}
This is the scalar equation represented by row \(i\) of \(A\) and \(B\).

\stepheading{why boundary values appear on the right-hand side}
For the first interior row, the lower-boundary terms at the two time levels are
\[
  a_1U_0^{n+1},
  \qquad
  a_1U_0^n.
\]
Moving the new-time boundary term to the right gives the combined contribution
\[
  \frac{\Delta\tau}{2}a_1
  \left(U_0^n+U_0^{n+1}\right).
\]
Similarly, the last row receives
\[
  c_{M-1}U_M^{n+1},
  \qquad
  c_{M-1}U_M^n,
\]
which combine as
\[
  \frac{\Delta\tau}{2}c_{M-1}
  \left(U_M^n+U_M^{n+1}\right).
\]
Hence
\[
  g^{n,n+1}
  =
  \frac{\Delta\tau}{2}
  \begin{bmatrix}
    a_1(U_0^n+U_0^{n+1})\\
    0\\
    \vdots\\
    0\\
    c_{M-1}(U_M^n+U_M^{n+1})
  \end{bmatrix}.
\]

\begin{implementationcheck}
This matches \path{src/solvers/cn_psor.py}: the first right-hand-side entry is
increased by
\(\tfrac12\Delta\tau\,a_1(U_0^n+U_0^{n+1})\), and the last by
\(\tfrac12\Delta\tau\,c_{M-1}(U_M^n+U_M^{n+1})\).
\end{implementationcheck}

\subsubsection{Boundary conditions used in the project}
For an American put,
\[
  U(0,\tau)=K,
  \qquad
  U(S_{\max},\tau)=0.
\]
At zero spot, immediate exercise returns \(K\). At a sufficiently large spot,
the put value approaches zero.

For a dividend-paying American call,
\[
  U(0,\tau)=0,
\]
\[
  U(S_{\max},\tau)
  =
  \max\!\left(
    S_{\max}-K,\,
    S_{\max}e^{-q\tau}-Ke^{-r\tau},\,
    0
  \right).
\]
The upper boundary retains both immediate-exercise value and the discounted
European far-field expression.

For European validation with projection disabled, the put uses
\[
  U(0,\tau)=Ke^{-r\tau},
  \qquad
  U(S_{\max},\tau)=0,
\]
and the call uses
\[
  U(0,\tau)=0,
  \qquad
  U(S_{\max},\tau)=S_{\max}e^{-q\tau}-Ke^{-r\tau}.
\]

\stepheading{the American linear complementarity problem}
Let
\[
  d^n=BU^n+g^{n,n+1}.
\]
A European step solves the equality \(AU^{n+1}=d^n\). For an American option,
that equality can produce a value below the payoff. The discrete obstacle
problem is instead
\[
  AU^{n+1}\geq d^n,
  \qquad
  U^{n+1}\geq\payoff,
\]
\[
  \left(U^{n+1}-\payoff\right)^\top
  \left(AU^{n+1}-d^n\right)=0.
\]
The first inequality requires a nonnegative discrete PDE residual. The second
enforces the payoff obstacle. The final equality imposes complementarity:
each node is either continuing with zero PDE residual or exercising with zero
premium.

\begin{commonconfusion}
An American option is not priced by solving the European system and then
reporting whichever nodes happen to lie above the payoff. The obstacle is part
of the system itself, so the solution method must enforce it during the solve.
\end{commonconfusion}

\stepheading{projected successive over-relaxation}
At one time step, write \(Ax=d\) for the interior system and let
\(\payoff_i\) be the obstacle at node \(i\). Starting from an old iterate
\(x^{(k)}\), a Gauss--Seidel candidate is
\[
  y_i^{(k+1)}
  =
  \frac{
    d_i
    -\sum_{j<i}A_{ij}x_j^{(k+1)}
    -\sum_{j>i}A_{ij}x_j^{(k)}
  }{A_{ii}}.
\]
The first sum uses already updated entries; the second uses entries from the
previous sweep. Relaxation gives
\[
  x_i^{\mathrm{relaxed}}
  =
  x_i^{(k)}
  +\omega\left(y_i^{(k+1)}-x_i^{(k)}\right).
\]
Projection then enforces the obstacle:
\[
  x_i^{\mathrm{new}}=\max\!\left(\payoff_i,
  x_i^{\mathrm{relaxed}}\right).
\]
The parameter \(\omega\in(0,2)\) controls relaxation. Values above one may
accelerate convergence for suitable systems, but excessive relaxation can
destabilize the iteration.

\subsubsection{Stopping and complementarity diagnostics}
The implemented iteration stops when
\[
  \lVert x^{(k+1)}-x^{(k)}\rVert_\infty<\text{tol},
\]
or reports failure after the maximum number of sweeps. A small update is
necessary but is not the full financial diagnostic. After solving, define
\[
  z=Ax-d.
\]
The LCP checks are
\[
  x-\payoff\geq0,
  \qquad
  z\geq0,
  \qquad
  (x-\payoff)^\top z\approx0.
\]
These separately test obstacle feasibility, residual feasibility, and
complementarity.

\begin{readerintuition}
PSOR alternates two actions: move toward the linear-system solution, then clip
the result back to the financially admissible side of the payoff. Convergence
is credible only when both the numerical update and the LCP residual checks are
small.
\end{readerintuition}

\subsection{Implementation Details}
The tridiagonal operator is exposed by
\path{src/solvers/cn_psor.py} for sign and matrix tests. Boundary formulas live
in \path{src/solvers/boundary_conditions.py}. The code uses the payoff as the
initial condition at \(\tau=0\), forms \(A\) and \(B\) once per solve, adjusts
the first and last right-hand-side entries at every time step, and applies PSOR
only when American projection is enabled.

\subsection{Validation and Diagnostics}
The existing solver audit tests verify the operator on a linear function, the
signs in \(A\) and \(B\), put and dividend-call boundary values, and PSOR
projection/complementarity. European closed-form and no-dividend call controls
provide independent checks beyond internal matrix consistency.

\subsection{Results}
The verified smoke-grid maximum error is \(2.782\times10^{-4}\) for both the
European put validation and the no-dividend American call control. The American
put obstacle violation is reported as zero, with complementarity error
\(4.009\times10^{-12}\). These are engineering-scale validation results, not a
claim of universal or publication-grade accuracy.

\subsection{Failure Modes and Limitations}
Uniform grids resolve the payoff kink and free boundary only approximately.
Finite-domain truncation can contaminate the interior if \(S_{\max}\) is too
small. Crank--Nicolson begins from a nonsmooth payoff, and Rannacher smoothing
is not active in the current solver. These limitations are particularly
important for Gamma near maturity.

\subsection{Conclusion and Next Step}
The PDE has become a boundary-adjusted LCP at each time step. Later stages
extract the exercise boundary and Greeks, study resolution and runtime, and use
only diagnostic-passing solver outputs for machine-learning labels.
